\chapter{La economía soviética}
\label{ch:sovieteconomy}

\section{La discusión sobre los incentivos, y cómo se zanjó}

Entre 1963 y 1965 la dirigencia cubana sostuvo, en público y por escrito, uno de
los debates más serios de la historia de la economía socialista. De un lado, el
Che Guevara, entonces Ministro de Industrias, defendía un sistema presupuestario
de financiamiento en el que las empresas eran departamentos de una única firma
nacional sin contabilidad propia, las transferencias entre ellas no eran ventas y
a los trabajadores los movían incentivos morales ---la emulación, el trabajo
voluntario, la formación de un hombre nuevo--- antes que la paga. Del otro,
Carlos Rafael Rodríguez y los economistas del viejo partido comunista defendían
el autofinanciamiento al modo soviético de la época, en el que las empresas
llevaban libros, enfrentaban precios y pagaban primas.

El debate se condujo con más honestidad intelectual de la que suele recibir el
tema, y vale la pena leer la parte de Guevara; su objeción a usar categorías
capitalistas para construir el socialismo es coherente, y su advertencia de que
los incentivos materiales reproducirían la mentalidad a la que debían servir
resultó, en el caso soviético, profética. También se equivocaba en todo lo
operativo. Un sistema sin precios no puede saber cuál de dos procesos es más
barato, un sistema sin contabilidad por empresa no puede localizar una pérdida, y
un sistema en el que el esfuerzo se premia con un diploma produce, pasado el
primer entusiasmo, exactamente el esfuerzo que vale un diploma.

El debate no se zanjó por argumento. Guevara dejó el gobierno en 1965; su enfoque
se aplicó al máximo durante los últimos años sesenta, culminando en la abolición
del pequeño negocio en 1968 y en una economía donde el dinero casi había dejado
de funcionar ---alquileres, servicios, llamadas locales, entierros,
espectáculos deportivos y comedores escolares eran gratuitos, los salarios se
pagaban pero había progresivamente menos que comprar con ellos, y el ausentismo
alcanzó niveles que el Estado no podía controlar---. Después, a partir de 1970,
todo el aparato se sustituyó en bloque por el sistema soviético que Rodríguez
había defendido.

\section{Diez millones de toneladas}

El eje del giro fue una zafra.

Tras haber fracasado la industrialización acelerada de comienzos de los sesenta
---un intento de sustituir importaciones en todos los frentes, que naufragó por
la ausencia de insumos, repuestos y de nadie que hubiera dirigido una fábrica---
Cuba volvió al azúcar como fuente de acumulación, y en 1968 la dirigencia se
comprometió públicamente con una zafra de diez millones de toneladas en 1970,
cifra que superaba en cerca de un tercio cualquier cosa que el país hubiera
producido jamás. El compromiso se hizo en el registro característico de Castro:
como cuestión de honor nacional, con fecha, en un discurso.

La movilización para cumplirlo fue total. Las fábricas cerraron y sus obreros
fueron a la caña. Estudiantes, soldados, oficinistas y voluntarios extranjeros
fueron a la caña. El transporte, el combustible, los repuestos y la atención de
los cuadros se sacaron de todos los demás sectores. La zafra se extendió hasta
julio, arruinando la siembra de la temporada siguiente. El país produjo 8,5
millones de toneladas, la mayor zafra de la historia cubana antes o después, y el
esfuerzo rompió el resto de la economía: la producción cayó en la manufactura, el
transporte y la construcción, los cañaverales quedaron dañados por años, y los
centrales se habían operado más allá del mantenimiento.

El 26 de julio de 1970 Castro pronunció un discurso en el que admitía el fracaso
en términos inusualmente directos, asumía la responsabilidad personal y ofrecía
renunciar. La multitud se negó, como todos en la plaza y en la tribuna sabían que
haría. Lo que siguió al discurso, sin embargo, fue un cambio de sistema real, y
fue en dirección contraria a la que el discurso sugería. Cuba dejó de intentar
inventar su propio socialismo y adoptó el soviético.

\section{Institucionalización, 1972--1976}

Cuba ingresó al Consejo de Ayuda Mutua Económica, el CAME, en julio de 1972.
Adoptó planes quinquenales, una junta central de planificación y, en 1976, el
SDPE, el Sistema de Dirección y Planificación de la Economía, copia cercana del
modelo empresarial soviético con contabilidad de costos, indicadores de
rentabilidad y primas. El Primer Congreso del Partido Comunista se reunió en
1975, dieciséis años después de la revolución. Una nueva constitución se aprobó
por referendo en 1976, sustituyendo a la de 1940 tal como había sido enmendada;
definía al partido como fuerza dirigente de la sociedad y del Estado, y creaba el
Poder Popular, un sistema de asambleas municipales, provinciales y nacional
electas desde la base, con elección directa, secreta y competitiva a nivel
municipal y selección indirecta por encima. Las provincias se redibujaron de seis
a catorce. Siguieron un código civil y penal, y un código de familia.

El período 1971--1985 fue, en términos económicos corrientes, el mejor que ha
tenido Cuba. El crecimiento fue sólido, el consumo mejoró, el programa de
vivienda entregó a través de las microbrigadas ---equipos destacados de sus
empresas para construir bloques de apartamentos, con los pisos asignados de vuelta
a la empresa--- y los indicadores sociales del capítulo~\ref{ch:socialstate}
alcanzaron su máximo. Los cubanos nostálgicos de la revolución lo son, casi
siempre, de estos catorce años en particular.

\section{Quién estaba pagando de verdad}

La aritmética que hay debajo es lo más importante que hay que entender sobre
Cuba, porque todo lo que ocurre en la Parte II es consecuencia de su desaparición.

El comercio de Cuba con la Unión Soviética no era comercio a precios mundiales.
Bajo los protocolos azucareros de largo plazo, la URSS compraba el azúcar cubano
a un precio fijado muy por encima del mercado mundial y suministraba crudo a un
precio fijado muy por debajo. El níquel y los cítricos tenían arreglos parecidos.
El pago corría por una cuenta bilateral de compensación en rublos transferibles y
no en moneda convertible, de modo que el déficit persistente de Cuba se acumulaba
como una deuda que jamás había que saldar en nada que la Unión Soviética pudiera
gastar en otra parte.

Encima de las condiciones de precio había un segundo mecanismo, y es el que suele
omitirse. A Cuba se le asignaba más petróleo soviético del que consumía y se le
permitía reexportar el excedente por divisa convertible. A mediados de los ochenta
esa reexportación era una de las mayores fuentes individuales de divisas de Cuba
---en algunos años la mayor, por delante del azúcar---.\unv{} A Cuba se le estaba
entregando, en efecto, una renta petrolera por parte de un país que tenía petróleo
y quería un aliado en el Caribe.

Las magnitudes no son pequeñas y no están seriamente en disputa. En 1987 la Unión
Soviética pagó a Cuba el equivalente de unos 42 centavos de dólar por libra de
azúcar frente a un precio mundial inferior a 7 centavos ---más de seis veces el
mercado---, y una evaluación de la CIA de comienzos de la década encontró una
proporción comparable sobre un volumen de 3,2 millones de toneladas. Convertir la
relación entera a dólares es más difícil, porque depende de qué tasa se aplique a
los rublos transferibles, que es el problema que describía la nota sobre las
fuentes. Pero las estimaciones habituales sitúan el flujo combinado en torno al 23
por ciento del producto cubano entre 1985 y 1988, llegando a algo más de cuatro
mil millones de dólares en un solo año \citep{mesalago2000}.\unv{}
Aproximadamente una cuarta parte del ingreso nacional, de un solo patrón,
sostenida durante más de una década.

Junto a ello, Cuba se había endeudado en divisa convertible con bancos y
gobiernos occidentales en los años setenta, cuando los precios del azúcar
estuvieron brevemente altos y el crédito cubano era bueno. Cuando cayeron, no
pudo pagar: Cuba suspendió los pagos de su deuda en moneda convertible en 1986 y
está esencialmente en cesación de pagos con el Club de París desde entonces. Ese
impago tiene treinta y nueve años en 2026 y es la razón de que Cuba no pueda
endeudarse, que es la razón de que no pueda invertir, restricción que el
capítulo~\ref{ch:money} tiene que resolver antes de que nada más en la Parte III
sea financiable.

\section{El fracaso que importaba: ninguna diversificación}

\begin{figure}[t]
\centering
\includegraphics[width=\linewidth]{figs/comecon}
\caption{El subsidio, y lo que compró. (a)~El mecanismo: en 1987 la Unión
Soviética pagó más de seis veces el precio mundial por el azúcar cubano, y
suministró petróleo por debajo de él. (b)~El resultado: la participación del
azúcar en las exportaciones cubanas de mercancías apenas se movió a lo largo de
treinta años de planificación, la tasa de inversión más alta de América Latina y
una membresía en el CAME organizada en torno a la especialización internacional
---y después se derrumbó cuando el comprador desapareció---. Un precio garantizado
elimina la única señal que le dice a un productor que llegue a ser bueno en otra
cosa. El panel~(a) son dos valores medidos y no una serie, porque los datos de
precio año por año no son públicos de forma fiable y trazar una línea entre dos
puntos sería inventar la forma; los valores posteriores a 1989 en~(b) son
estimaciones y llevan la marca \unv{}. Generada por \texttt{figs/comecon.py}.}
\label{fig:comecon}
\end{figure}

He aquí el hecho que condena al período entero, y la Fig.~\ref{fig:comecon}(b) es
su retrato.

En 1958 el azúcar era alrededor de cuatro quintas partes de las exportaciones
cubanas. En 1989, después de treinta años de planificación, tres décadas con las
tasas de inversión más altas de América Latina, una membresía en el CAME
explícitamente organizada en torno a la especialización internacional y la fuerza
de trabajo mejor educada de la región, el azúcar seguía siendo alrededor de tres
cuartas partes de las exportaciones cubanas \citep{onei,mesalago2000}.\unv{}

Todo lo demás era pequeño. El níquel, de los yacimientos de Moa y Nicaro, era el
segundo. Seguían el tabaco, los cítricos, la pesca y el ron. La base industrial
levantada en los setenta era ensamblaje y sustitución de importaciones para el
mercado interno, funcionaba con insumos soviéticos y no podía exportar fuera del
CAME porque sus costos y su calidad los fijaba un mercado que no ponía a prueba
ninguna de las dos cosas. Se inició a comienzos de los ochenta una capacidad
farmacéutica y biotecnológica ---sobre la que hay mucho más en el
capítulo~\ref{ch:life}, porque es la única excepción y el único activo---, pero
en 1989 estaba a una década de tener productos.

La razón del fracaso no es misteriosa y no es el embargo. Un mercado protegido a
precios garantizados elimina la única señal que le dice a un productor que llegue
a ser bueno en otra cosa. Cuba tuvo, durante dieciocho años, los términos de
intercambio más dulces que se le hayan ofrecido jamás a un país pequeño, y los
usó para seguir siendo una economía azucarera con un muy buen sistema escolar.
Esa es la lección económica central de la Parte I, y el
capítulo~\ref{ch:premises} la enuncia como restricción del diseño: \emph{un
comprador garantizado a un precio político no es un mercado, y una economía
construida sobre uno no se desarrolla.} Se aplicaba a la cuota azucarera
norteamericana antes de 1959 y se aplicó a los protocolos soviéticos después, y el
capítulo~\ref{ch:venezuela} lo mostrará aplicándose una tercera vez al petróleo
venezolano.

\section{La Rectificación, 1986}

En 1980, presionado por la escasez de alimentos, Cuba abrió mercados libres
campesinos donde las cooperativas y los agricultores privados podían vender el
excedente sobre su cuota estatal al precio que el mercado soportara. Funcionaron:
apareció comida. Produjeron también una desigualdad visible ---intermediarios,
algunos campesinos a los que les iba muy bien, precios que los más pobres no
podían pagar--- y quejas por especulación.

En 1986 Castro lanzó la Rectificación de Errores y Tendencias Negativas. Cerró
los mercados campesinos. Prohibió la construcción privada de viviendas para la
venta y buena parte del trabajo artesanal privado. Recortó el sistema de primas
empresariales del SDPE, denunció los incentivos materiales, revivió el trabajo
voluntario y los contingentes, y rehabilitó explícitamente y por su nombre el
pensamiento económico de Guevara.

Lo decisivo es el momento. En 1986 la Unión Soviética iniciaba la perestroika.
Cuba respondió al giro de su patrón hacia los mercados girando en sentido
contrario, y lo hizo dieciocho meses antes de que empezaran a desaparecer los
términos de intercambio que hacían asequible esa decisión. La Rectificación es el
ejemplo más nítido del patrón que este libro nombrará en el
capítulo~\ref{ch:ledgerch} y que se repite a lo largo de la Parte II: \emph{Cuba
liberaliza bajo presión y recentraliza en cuanto la presión cede, y la
recentralización llega siempre justo antes del momento en que es menos
asequible.}

\begin{ledger}{la economía soviética, 1963--1989}
\emph{Logrado.} Catorce años de crecimiento sólido y el nivel de vida más alto
que ha producido la revolución. Un Estado social universal financiado hasta el
final. Vivienda masiva mediante las microbrigadas. Pleno empleo, en el sentido de
que todo el mundo tenía un puesto. Formas institucionales ---una constitución, un
congreso del partido, asambleas municipales electas, derecho codificado--- donde
antes había habido improvisación por discurso. Un sistema educativo cuyo producto
era, y sigue siendo, el mejor activo del país. La fundación de una capacidad
biotecnológica que sobreviviría a todo lo demás aquí enumerado.

\emph{Costo.} Treinta años de acumulación que dejaron la canasta exportadora
donde estaba en 1958, tres cuartas partes azúcar, porque los precios garantizados
eliminaron la razón para cambiarla. Una economía estructuralmente dependiente de
transferencias de entre una quinta y una tercera parte del ingreso nacional, de un
solo patrón, sin plan alguno para su terminación. Un impago en divisa en 1986 que
ha dejado al país fuera de los mercados de crédito desde entonces. Un sistema
empresarial incapaz de calcular, y una fuerza de trabajo cuya paga se había
desacoplado de lo que hacía. Y una decisión en 1986 de revertir la única reforma
que visiblemente había funcionado, tomada en el último momento en que revertirla
era todavía sobrevivible.
\end{ledger}
